\section{Neural-Guided Semantic Proxy (NGSP): Privacy-Preserving
Mediation of Clinical Trial Documents to Consumer LLM
Interfaces}\label{neural-guided-semantic-proxy-ngsp-privacy-preserving-mediation-of-clinical-trial-documents-to-consumer-llm-interfaces}

\textbf{Abstract.} Clinical trial staff routinely paste sensitive
documents --- Serious Adverse Event narratives, protocol excerpts,
monitoring reports --- into consumer LLM chat interfaces. Existing
defenses either over-redact (destroying utility) or under-redact
(leaking quasi-identifiers and protocol structure). We present NGSP, a
Neural-Guided Semantic Proxy that runs Gemma 4 locally as a privacy
filter composed with deterministic HIPAA Safe Harbor stripping and a
calibrated (ε, δ)-differentially-private bottleneck. A three-way router
directs inputs to one of three paths: abstract query synthesis (breaking
injective entity linkage), DP-noised proxy generation (providing formal
guarantees), or local-only answering (for content-inseparable tasks). We
evaluate the system against five adversarial attack classes --- verbatim
scan, cross-encoder similarity, trained inversion (primary threat
model), membership inference, and utility regression --- on a synthetic
clinical trial corpus of 1,200 documents with ground-truth sensitive
span annotations. At ε = 3.0, δ = 1e-5, we report {[}results pending
execution of experiments{]}. The hypothesis that the composed system
simultaneously bounds inversion span-recovery F1 ≤ 0.09 and preserves
utility ≥ 0.85 is {[}pending{]}. We report all results including
negative findings.

\begin{center}\rule{0.5\linewidth}{0.5pt}\end{center}

\subsection{1. Introduction}\label{introduction}

The gap between sanctioned clinical data handling systems and the
practical utility of consumer LLM interfaces creates a structural
pressure for policy violation. A medical writer preparing an SAE
narrative for submission, a CRA summarizing a monitoring report, or a
biostatistician querying an analysis plan will reach for a faster,
better-quality tool when their sanctioned alternatives are slower or
produce inferior output. Data Loss Prevention (DLP) systems respond by
blocking pasting entirely, at the cost of driving users to personal
devices. Business Associate Agreement (BAA) wrappers add legal coverage
but do not reduce semantic information content in what is transmitted.

NGSP takes a different approach: rather than blocking or permitting, it
transforms. The local Gemma model acts as a privacy-aware intermediary
that rewrites, abstracts, or adds calibrated noise before any text
leaves the local device, allowing the high-quality cloud LLM to answer a
privacy-safe variant of the original question.

\textbf{Contributions:} 1. A three-path neural router that classifies
queries into abstract-extractable, DP-tolerant, and local-only
categories based on the separability of task intent from sensitive
content. 2. A session-scoped Rényi DP accounting scheme for per-user
privacy budget tracking with hard budget enforcement. 3. A five-attack
adversarial test harness with ground-truth sensitive span annotations on
a 1,200-document synthetic corpus. 4. A quantitative evaluation of the
privacy-utility tradeoff across six ε values, with an ablation study
isolating each defense component's contribution.

\subsection{2. Related Work}\label{related-work}

\textbf{Differential privacy for text.} Feyisetan et al.~(2020) apply
local DP to word embeddings for text privatization, and Yue et
al.~(2021) propose a natural-text sanitization mechanism. More recent
work has sharpened both the mechanism design and the critique surface:
Mattern et al.~(2022) analyze the limits of word-level DP, Chen et
al.~(2023) introduce customized token-level sanitization, and Awon et
al.~(2025), Zhang et al.~(2025), and Tong et al.~(2025) extend the
modern text-sanitization literature with semantically-aware and
attack-aware formulations. Our approach differs in that we apply DP to
hidden-state activations rather than token distributions, and we compose
DP with a non-injective query synthesis path that provides stronger
practical privacy for the dominant query class.

\textbf{LLM privacy and leakage.} Carlini et al.~(2021) demonstrate
training-data extraction from large language models. Mireshghallah et
al.~(2022) quantify privacy risks via membership inference attacks, and
Huang et al.~(2022) study leakage of personal information from large
pretrained language models. Our threat model is more specific: we target
quasi-identifier leakage in proxy text, not only model memorization at
training time.

\textbf{Clinical NLP de-identification.} Lample et al.~(2016),
Dernoncourt et al.~(2017), Liu et al.~(2017), and Friedrich et
al.~(2019) anchor the transition from sequence-labeling-based PHI
removal toward privacy-preserving text representations for medical
records. Standard de-identification systems focus on the 18 HIPAA Safe
Harbor identifiers using rule-based and NER approaches. NGSP's Safe
Harbor component applies the same approach but treats it as the
\emph{first} layer of a composed defense, not the only one.

\textbf{Privacy accounting and practical deployment.} Mironov (2017)
introduces Rényi differential privacy, and Wang et al.~(2019) provide
the subsampled analytical moments accountant used widely in practical
composition analyses. NGSP adopts this accounting perspective, but
couples it to a local-model mediation architecture and a task-aware
router designed for interactive drafting workflows on commodity
hardware.

\subsection{3. Threat Model}\label{threat-model}

\textbf{System model.} A user on a corporate workstation runs an NGSP
browser extension that intercepts text pasted into a consumer LLM chat
interface. The NGSP process runs locally and holds the Gemma model
weights. The extension communicates with the cloud LLM (Anthropic
Claude) via the NGSP pipeline, forwarding only the privacy-processed
proxy text.

\textbf{Adversary model.} A passive adversary with access to: (a) all
proxy texts forwarded to the cloud LLM, (b) a background corpus of
clinical documents from the same distribution, (c) unlimited computation
for offline attack model training. The adversary cannot observe: local
Gemma activations, the entity\_map substitution dictionary, DP noise
seeds, or routing decisions.

\textbf{Attack goals:} 1. \emph{Verbatim recovery}: Does any sensitive
span value appear literally in the proxy? 2. \emph{Semantic proximity}:
Is the proxy semantically close enough to the original that
retrieval-based re-identification is feasible? 3. \emph{Span
prediction}: Given a proxy, can a trained model predict which tokens
correspond to sensitive spans in the original? 4. \emph{Membership
inference}: Can an adversary determine whether a specific entity
(compound code) was mentioned in the original document? 5. \emph{Utility
preservation} (inverted): Is the proxy useful enough to answer the
original task?

\subsection{4. System Description}\label{system-description}

\subsubsection{4.1 Safe Harbor Stripper}\label{safe-harbor-stripper}

Deterministic regex patterns for 14 of the 18 HIPAA Safe Harbor
identifiers (dates, phone numbers, SSN, MRN, ZIP codes, email addresses,
URLs, IP addresses, device IDs, account numbers, certificate numbers,
vehicle IDs, web URLs, biometric IDs). A Gemma NER pass handles the
remaining four (names, geographic subdivisions, full-face photographs as
described text, and free-text identifiers). Detected entities are
replaced with typed placeholders
(\texttt{\textless{}PERSON\_1\textgreater{}},
\texttt{\textless{}DATE\_1\textgreater{}}) and stored in a local
\texttt{entity\_map}. The stripped text and entity\_map are the inputs
to the router.

\subsubsection{4.2 Router}\label{router}

Gemma 4 receives the stripped text and classifies it as one of three
paths. The classification prompt instructs Gemma to assess whether the
\emph{task intent} is separable from \emph{entity-specific content}:

\begin{itemize}
\tightlist
\item
  \texttt{abstract\_extractable}: The user's question can be rephrased
  without reference to any entity in the document. Example: ``What are
  the standard criteria for a Grade 3 AE?'' --- the answer is
  independent of which patient or compound is involved.
\item
  \texttt{dp\_tolerant}: The answer requires conveying the document's
  content, but semantic fidelity can tolerate distortion. Example:
  ``Summarize this SAE narrative'' --- a noised proxy preserves enough
  structure for a useful summary.
\item
  \texttt{local\_only}: Task and content are inseparable. Example: ``Is
  this compound code consistent with the protocol amendment?'' ---
  answering requires the specific entity values, which must never leave
  the device.
\end{itemize}

The router falls back to \texttt{dp\_tolerant} on parse errors or model
failures.

\subsubsection{4.3 Query Synthesizer (Abstract-Extractable
Path)}\label{query-synthesizer-abstract-extractable-path}

For abstract-extractable inputs, Gemma generates a new self-contained
question that captures the user's task intent without mentioning any
entity from the stripped input. The synthesized query is sent to the
cloud LLM. Because the mapping is non-injective (many inputs can produce
the same synthesized query), this path provides strong practical privacy
even without formal DP guarantees: there is no function from proxy →
original entity set.

\subsubsection{4.4 DP Bottleneck (DP-Tolerant
Path)}\label{dp-bottleneck-dp-tolerant-path}

Hidden states at the last layer of Gemma's generation pass are
mean-pooled over new tokens to produce a single vector
\texttt{h\ ∈\ ℝ\^{}d}. The mechanism:

\begin{enumerate}
\def\labelenumi{\arabic{enumi}.}
\tightlist
\item
  \textbf{Clip}: \texttt{h\_clip\ =\ h\ ·\ min(1,\ C/‖h‖₂)}, C = 1.0
  (bounded L2 sensitivity)
\item
  \textbf{Noise}: \texttt{h\_noisy\ =\ h\_clip\ +\ N(0,\ σ²I)}, σ =
  3.298/ε
\item
  \textbf{Decode}: Gemma paraphrase conditioned on the noisy embedding
  as a semantic hint
\item
  \textbf{Account}: Rényi DP accountant accumulates R\_α = α/(2σ²) per
  call; converts to (ε, δ)-DP at query time
\end{enumerate}

Per-session budget: ε\_cap = 3.0, δ = 1e-5.
\texttt{BudgetExhaustedError} is raised when the cumulative session ε
exceeds the cap.

\subsubsection{4.5 Answer Applier}\label{answer-applier}

The cloud LLM response is returned to the local process, where the
entity\_map is re-applied (longest-key-first to avoid partial match
collisions: \texttt{\textless{}PERSON\_12\textgreater{}} is replaced
before \texttt{\textless{}PERSON\_1\textgreater{}}). The final answer
with original entity values restored is returned to the user.

\subsection{5. Experimental Setup}\label{experimental-setup}

\textbf{Corpus.} 1,200 synthetic clinical trial documents (500 SAE
narratives, 200 protocol excerpts, 200 monitoring reports, 300 CSR draft
paragraphs). Generated with seed 42; reproducible via
\texttt{src/data/}. Ground-truth sensitive span annotations provided by
the \texttt{annotate()} oracle.

\textbf{Evaluation split.} Inversion attacker: 80\% train / 20\% eval.
Membership inference: positive/negative split by entity presence.
Utility: 50-document random sample (seed 42).

\textbf{Hardware.} Apple Silicon (MPS) for Gemma 4 inference and
DistilBERT fine-tuning. CPU fallback tested. \texttt{device\_map="mps"}
avoided (known accelerate pre-allocation issue on MPS; model loaded
without device\_map, then \texttt{.to("mps")}).

\textbf{Models.} Gemma 4 E2B-it (\texttt{google/gemma-4-E2B-it}),
DistilBERT base uncased, \texttt{all-MiniLM-L6-v2}
(sentence-transformers).

\subsection{6. Results}\label{results}

\emph{{[}Populate from
\texttt{experiments/results/attacks\_eps3.0.json},
\texttt{experiments/results/ablations.json},
\texttt{experiments/results/calibration.json} after running
experiments.{]}}

See \texttt{paper/results.md} for the full results tables.

\subsubsection{6.1 Privacy-Utility
Tradeoff}\label{privacy-utility-tradeoff}

{[}Insert Figure: \texttt{paper/figures/epsilon\_utility\_curve.png}{]}

The calibration sweep over ε ∈ \{0.5, 1.0, 2.0, 3.0, 5.0, 10.0\} reveals
the operating envelope. {[}Populate with actual findings from
calibration.json.{]}

\subsubsection{6.2 Attack Results at ε =
3.0}\label{attack-results-at-ux3b5-3.0}

\textbf{Attack 1 (Verbatim):} {[}Values from
\texttt{results{[}"verbatim"{]}} in attacks\_eps3.0.json.{]}

\textbf{Attack 2 (Similarity):} {[}Values from
\texttt{results{[}"similarity"{]}}.{]}

\textbf{Attack 3 (Inversion):} {[}Values from
\texttt{results{[}"inversion"{]}}.{]} - H₁ \textbf{{[}PASS/FAIL{]}}:
overall\_f1 = {[}value{]} vs.~threshold 0.09

\textbf{Attack 4 (Membership):} {[}Values from
\texttt{results{[}"membership"{]}}.{]}

\textbf{Attack 5 (Utility):} {[}Values from
\texttt{results{[}"utility"{]}}.{]} - H₂ \textbf{{[}PASS/FAIL{]}}:
utility\_ratio = {[}value{]} vs.~threshold 0.85

\subsubsection{6.3 Ablation Study}\label{ablation-study}

{[}Insert Table from \texttt{experiments/results/ablations.json}.{]}

The ablation isolates each component's contribution. Key question for
H₃: does any single component achieve both H₁ and H₂?

\subsection{7. Discussion}\label{discussion}

See \texttt{paper/discussion.md} for the full discussion. Key points:

\textbf{Privacy-utility curve.} {[}Characterize after results.{]}

\textbf{Proxy decoder limitation.} The v1 proxy decoder uses the noisy
hidden state as a soft hint to a Gemma paraphrase call rather than as a
hard decoding constraint. This provides approximate rather than strict
DP guarantees for the proxy text. A dedicated decoder network would be
the correct next step.

\textbf{Synthetic corpus scope.} All evaluation is on synthetic data.
Generalization to real clinical documents requires a separately governed
study.

\subsection{8. Limitations}\label{limitations}

\begin{enumerate}
\def\labelenumi{\arabic{enumi}.}
\tightlist
\item
  Synthetic corpus only; generalization to real data is untested.
\item
  Passive adversary only; adaptive adversaries are out of scope.
\item
  DP path proxy decoder provides approximate, not strict, guarantees.
\item
  Rare-entity membership inference (\textless{} 5 appearances) is not
  evaluated.
\item
  Latency impact on user adoption is not measured.
\item
  Gemma 4 at 2B parameters; larger models may produce better routing and
  proxy quality.
\end{enumerate}

\subsection{9. Conclusion}\label{conclusion}

NGSP demonstrates the feasibility of a three-layer composed privacy
architecture (Safe Harbor + neural routing + DP bottleneck) as a
practical alternative to DLP blocking for clinical trial document
handling. The five-attack adversarial evaluation provides a quantitative
characterization of the privacy-utility tradeoff at multiple ε values.
{[}Conclude with findings after experiments run.{]}

\begin{center}\rule{0.5\linewidth}{0.5pt}\end{center}

\subsection{Appendix A --- Repository
Structure}\label{appendix-a-repository-structure}

\begin{verbatim}
ngsp-clinical/
├── src/ngsp/         # Pipeline implementation
├── src/attacks/      # Five attack classes
├── src/data/         # Synthetic corpus generators
├── experiments/      # Attack battery and calibration runners
├── tests/            # Pytest test suite (74 tests, 0 failures)
├── paper/            # This document and supporting files
└── scripts/          # Setup, download, comment checker
\end{verbatim}

\subsection{Appendix B --- Running the
Experiments}\label{appendix-b-running-the-experiments}

\begin{Shaded}
\begin{Highlighting}[]
\CommentTok{\# Prerequisites}
\ExtensionTok{./scripts/setup.sh}
\ExtensionTok{python3}\NormalTok{ scripts/download\_gemma.py   }\CommentTok{\# requires HF\_TOKEN with Gemma 4 access}
\FunctionTok{cp}\NormalTok{ .env.example .env                }\CommentTok{\# fill in ANTHROPIC\_API\_KEY}

\CommentTok{\# Smoke tests (no model load, no API calls)}
\ExtensionTok{pytest} \AttributeTok{{-}q}\NormalTok{ tests/test\_attacks.py tests/test\_mock\_data.py}

\CommentTok{\# Full attack battery at ε = 3.0}
\ExtensionTok{python3}\NormalTok{ experiments/run\_attacks.py }\AttributeTok{{-}{-}epsilon}\NormalTok{ 3.0 }\AttributeTok{{-}{-}n{-}docs}\NormalTok{ 50}

\CommentTok{\# Ablation study}
\ExtensionTok{python3}\NormalTok{ experiments/ablations.py }\AttributeTok{{-}{-}n{-}docs}\NormalTok{ 20}

\CommentTok{\# Privacy{-}utility calibration sweep}
\ExtensionTok{python3}\NormalTok{ experiments/calibrate\_epsilon.py}

\CommentTok{\# Comment invariant check}
\ExtensionTok{python3}\NormalTok{ scripts/check\_comments.py }\AttributeTok{{-}{-}path}\NormalTok{ src/}
\end{Highlighting}
\end{Shaded}

\subsection{Appendix C --- Sensitive Category
Taxonomy}\label{appendix-c-sensitive-category-taxonomy}

The \texttt{SensitiveCategory} enum in \texttt{src/data/schemas.py}
covers:

\textbf{HIPAA Safe Harbor (18):} NAME, DATE, PHONE, FAX, EMAIL, SSN,
MRN, HEALTH\_PLAN\_NUMBER, ACCOUNT\_NUMBER, CERTIFICATE\_NUMBER,
VEHICLE\_ID, DEVICE\_ID, URL, IP\_ADDRESS, BIOMETRIC\_ID, PHOTO,
GEOGRAPHIC\_SUBDIVISION, OTHER\_UNIQUE

\textbf{Quasi-identifiers:} COMPOUND\_CODE, SITE\_ID, DOSE, INDICATION,
EFFICACY\_VALUE, AE\_GRADE, TIMING, COMORBIDITY

\textbf{MNPI:} INTERIM\_RESULT, AMENDMENT\_RATIONALE,
REGULATORY\_QUESTION
